\section{First-principles component models}
\label{sec:components}

\Cref{mod:sm,mod:fs} introduced two constants, $S_m$ and $f_s$, that carry a
great deal of weight and are not laws. This section replaces both with models
derived from geometry and materials, at the cost of losing the closed form. The
engine keeps both layers and reports the implied values of $S_m$, $f_s$ and
$\FM$ so that the lumped assumptions can be audited against the detailed model.

\subsection{Arms: a cantilevered thin-walled tube}

\begin{model}[Rotor arm]
\label{mod:arm}
Each of the $N$ arms is a thin-walled circular tube of mean radius $r$, wall
thickness $t\ll r$, length $L$, elastic modulus $E$, density $\rho_m$ and
allowable stress $\sigma_{\mathrm{all}}$, built in at the hub and carrying the
rotor as a point load at the tip.
\end{model}

\paragraph{Length.} Adjacent rotor disks must not overlap. With $N$ hubs on a
circle of radius $L$, the centre-to-centre spacing of adjacent hubs is
$2L\sin(\pi/N)$, so non-overlap requires $2L\sin(\pi/N)\ge 2R$, i.e.
\begin{equation}
  L = \frac{k_{\mathrm{gap}}\,R}{\sin(\pi/N)}, \qquad k_{\mathrm{gap}}\ge1 ,
  \label{eq:armlength}
\end{equation}
with $k_{\mathrm{gap}}$ a clearance factor. The tip-to-tip span is
$2(L+R)$ and the \emph{reach}, the distance from the vehicle centre to the
nearest blade tip, is
\begin{equation}
  \mathcal{R} \coloneqq L + R .
  \label{eq:reach}
\end{equation}
\Cref{eq:reach} is trivial and turns out to be one of the most consequential
quantities in the paper; see \cref{sec:results}.

\paragraph{Section properties.} For a thin-walled tube the second moment of
area about a diameter is
\begin{equation}
  I = \pi r^3 t + O(t^2), \qquad
  \text{mass per unit length } \mu = 2\pi r t \rho_m .
  \label{eq:tube}
\end{equation}

\paragraph{Strength.} Under the limit load $F = \mathrm{SF}\,\lambda m g/N$ at
the tip, the root bending moment is $M=FL$ and the extreme-fibre stress is
$\sigma = Mr/I$. Requiring $\sigma\le\sigma_{\mathrm{all}}$ gives
\begin{equation}
  t \ge t_{\sigma} = \frac{F L}{\pi r^2 \sigma_{\mathrm{all}}} .
  \label{eq:tstress}
\end{equation}

\paragraph{Stiffness.} Euler-Bernoulli theory gives the tip deflection of a
cantilever under a tip load as $\delta = FL^3/(3EI)$. Imposing
$\delta\le L/\Lambda$ for a slenderness limit $\Lambda$ (we use $\Lambda=100$)
gives
\begin{equation}
  t \ge t_{\delta} = \frac{\Lambda F L^{2}}{3\pi E r^{3}} .
  \label{eq:tstiff}
\end{equation}

\paragraph{Dynamics.} The fundamental bending frequency of a uniform cantilever
of mass per unit length $\mu$ is
\begin{equation}
  f_1 = \frac{1.875^2}{2\pi}\sqrt{\frac{EI}{\mu_{\mathrm{eff}} L^4}},
  \qquad
  \mu_{\mathrm{eff}} = \mu + \frac{3 m_{\mathrm{tip}}}{L} ,
  \label{eq:freq}
\end{equation}
where $1.875$ is the first root of $\cos\zeta\cosh\zeta+1=0$ and the correction
to $\mu$ is the Rayleigh lumped-mass approximation for a tip mass
$m_{\mathrm{tip}}$ (motor plus propeller). The wall thickness actually used is
\begin{equation}
  t = \max\{t_\sigma,\; t_\delta,\; t_{\min}\} ,
  \label{eq:wall}
\end{equation}
and the engine records which of the three governs, because the answer is
informative: for the designs of \cref{sec:results} it is stiffness or minimum
gauge, never strength, with a working stress of \SI{44}{\mega\pascal} against
an allowable of \SI{350}{\mega\pascal}.

\begin{remark}[Consistency with \cref{as:rigid}]
\Cref{eq:freq} is what makes the rigid-body assumption checkable. With
$f_1\approx\SI{32}{\hertz}$ and an attitude loop at
$\sqrt{K_R}/2\pi\approx\SI{1.0}{\hertz}$ (\cref{sec:control}) the separation is
a factor of $32$, comfortably in the regime where treating the airframe as rigid
does not corrupt the control design.
\end{remark}

\paragraph{Frame total.} Guards are rings around each disk, so their mass scales
with circumference, not area:
\begin{equation}
  m_{\mathrm{guard}} = k_g\, N\, 2\pi R .
  \label{eq:guards}
\end{equation}
The frame is then
\begin{equation}
  \mstr = N\mu L\,(1+k_{\mathrm{hub}}) + m_{\mathrm{guard}}
          + m_{\mathrm{boom}} + m_{\mathrm{fixed}} ,
  \label{eq:frame}
\end{equation}
with $k_{\mathrm{hub}}$ accounting for the centre plate, motor mounts and
fasteners, $m_{\mathrm{boom}}$ the optional display cantilever of
\cref{sec:results}, and $m_{\mathrm{fixed}}$ landing gear and wiring.

\subsection{Motors: sized by torque, with a scale effect}

The rotor demands a torque $Q_{\max} = \Pshaft^{\max}/\Omega_{\max}$ at the
sizing condition $T_{\max}=\lambda mg/N$ per rotor. A motor's continuous torque
is set by the shear stress its airgap can sustain over its rotor surface,
$Q \sim \tau_{\mathrm{airgap}}\,(\pi D_m L_m)\,(D_m/2)$, which for
geometrically similar machines gives $Q\propto D_m^3 \propto m_{\mathrm{mot}}$,
i.e.\ constant torque density. Real small outrunners do better than that as
they grow, because the ratio of cooling surface to loss volume improves and
because the winding fill factor rises. We therefore write

\begin{model}[Motor mass]
\label{mod:motor}
\begin{equation}
  Q = k_\tau \left(\frac{m_{\mathrm{mot}}}{m_{\mathrm{ref}}}\right)^{\!e}
      m_{\mathrm{mot}}
  \quad\Longleftrightarrow\quad
  m_{\mathrm{mot}}
  = \left(\frac{Q\, m_{\mathrm{ref}}^{e}}{k_\tau}\right)^{\!\frac{1}{1+e}} ,
  \label{eq:motormass}
\end{equation}
with $m_{\mathrm{ref}}=\SI{0.1}{\kilo\gram}$, $e\approx0.30$ and $k_\tau$
calibrated in \cref{sec:verification}.
\end{model}

Because the motor is sized by the burst condition $T_{\max}$, the peak rather
than continuous torque is the right basis. \Cref{eq:motormass} is sublinear in
$Q$, which matters: at $e=0.3$ a tenfold increase in torque costs a factor
$10^{1/1.3}\approx5.9$ in mass, not ten.

Speed controllers are sized by power, $m_{\mathrm{esc}} = \Pelec^{\max}/k_P$,
and propellers by an empirical areal law $m_{\mathrm{prop}} = k_p D^{2.6}
(N_b/2)$ fitted to carbon propellers.

\begin{remark}[The electrical chain that mass does not see]
\Cref{eq:motormass} determines the size of the motor but not whether it can be
driven. A rotor turning at $\Omega$ requires a back-EMF $\Omega/K_V$ below the
pack voltage with margin for the current-resistance drop, and the motor
constant satisfies $K_t\,[\si{\newton\metre\per\ampere}] = 9.5493/K_V$
$[\si{rpm\per\volt}]$. Sizing $K_V$ so that maximum rotor speed is reached at
the \emph{minimum} pack voltage, not the nominal or full one, is what keeps the
machine flyable at the bottom of a discharge. This constraint is decoupled from
the mass closure and is treated separately in the engine; it is what forces the
low-$K_V$, custom-wound motors of \cref{sec:results}.
\end{remark}

\subsection{Battery: two constraints, not one}

\Cref{mod:battery} sizes the pack by energy. A pack must also deliver power.
Writing $p_b$ for the usable specific power and $\varsigma$ for the state of
charge held in reserve,
\begin{equation}
  \mbat = \max\left\{
    \underbrace{\frac{E_{\mathrm{mission}}}{\eb\,(1-\varsigma)}}
      _{\text{energy}},\;
    \underbrace{\frac{\Pelec^{\max}}{p_b}}_{\text{power}}
  \right\} .
  \label{eq:batmax}
\end{equation}
Which term binds depends on the mission: for short sessions the power term
dominates and endurance is effectively free, while for long sessions the energy
term takes over. The engine reports which is active, because a design that is
power-limited responds to entirely different levers than one that is
energy-limited.

High-energy lithium-ion cells are not high-power cells. Sizing against the
burst condition $T_{\max}$ requires the burst rating; the continuous hover draw
must then be checked separately against the continuous rating. Both checks are
implemented.

\subsection{Cost of the first-principles layer}

Replacing \cref{mod:sm,mod:fs} by \cref{mod:arm,mod:motor} removes the closed
form: $\mstr(m)$ and $\mprop(m)$ are no longer proportional to $m$, so
\cref{eq:closure} is replaced by the implicit equation
\begin{equation}
  \mathcal{F}(m) \coloneqq m - \Big[\mfix + \mstr(m) + \mprop(m) + \mbat(m)\Big]
  = 0 ,
  \label{eq:fpclosure}
\end{equation}
solved numerically. For the designs studied $\mathcal{F}$ has the qualitative
shape of \cref{sec:fold}: negative for small $m$, positive on the branch where a
design exists, negative again beyond the fold. That shape is observed, not
proved. \Cref{lem:concave} is a statement about the lumped residual; it does
not transfer to $\mathcal{F}$, which contains clipped quantities, maxima
(\cref{eq:batmax}) and empirical fits. Near the fold the positive stretch can
also be very narrow, and \cref{sec:numerics} describes how the engine searches
for it and what it reports when it does not find it. A failed search is not a
proof that no mass closes.
